\title{DeepSeekMath: Pushing the Limits of Mathematical Reasoning in Open Language Models}

\begin{abstract}
Language models often struggle with mathematical reasoning because of its inherently intricate and organized framework. In this study, we present DeepSeekMath~7B, which extends the pre-training of DeepSeek-Coder-Base-v1.5 7B by incorporating an additional 120B tokens focused on mathematics, obtained from Common Crawl alongside a mix of natural language and programming code data. DeepSeekMath~7B achieves a notable score of 51.7\% on the challenging MATH benchmark at competition level, all without leveraging auxiliary toolkits or ensembling strategies, and performs comparably to advanced systems like Gemini-Ultra and GPT-4. When employing self-consistency across 64 generated samples, DeepSeekMath~7B attains 60.9\% accuracy on MATH. This improvement in mathematical reasoning performance is credited to two principal elements: first, the use of an extensively refined data filtering and selection process that unlocks the potential of open web data; and second, the introduction of Group Relative Policy Optimization (GRPO), an adaptation of Proximal Policy Optimization (PPO) that augments reasoning capabilities while optimizing PPO’s memory efficiency.
\end{abstract}

\section{Introduction}

Large language models (LLMs) have transformed the landscape of mathematical reasoning in AI, driving substantial progress in quantitative reasoning benchmarks~\citep{MATH} and geometry reasoning evaluations~\citep{trinh2024solving}. Additionally, these systems have shown substantial utility in aiding humans to solve intricate mathematical problems~\citep{tao}. Nonetheless, leading-edge models such as GPT-4~\citep{gpt4} and Gemini-Ultra~\citep{gemini} remain unavailable to the public, with the current generation of open-source models falling considerably short in terms of performance.

This paper presents DeepSeekMath, a language model tailored for mathematical reasoning that substantially surpasses other open-source models in mathematical tasks and nears the accuracy of GPT-4 on several academic evaluation sets. Central to this work is the construction of the DeepSeekMath Corpus, a pre-training dataset encompassing 120B math-specific tokens curated for quality at scale. To assemble this corpus, we leveraged a fastText-based classifier \citep{joulin2016fasttext} trained initially with positive examples drawn from OpenWebMath \citep{openwebmath} alongside a representative set of negative samples from various web domains. With the classifier, we iteratively identified new positive samples within the Common Crawl (CC), which underwent further manual curation to ensure relevance. The classifier was continually refined using these new annotations, enhancing its extraction accuracy. Our results demonstrate that this extensive corpus underpins significant advances: DeepSeekMath-Base 7B records a score of 64.2\% on GSM8K \citep{gsm8k} and 36.2\% on the MATH dataset \citep{MATH}, exceeding the performance of Minerva 540B \citep{minerva}. Owing to the corpus’s multilingual character, we also observe gains on Chinese mathematics benchmarks \citep{wei2023cmath,agieval}. We suggest that our methodologies for curating mathematical data offer a strong foundation for further study, and we anticipate considerable progress is still achievable.

DeepSeekMath-Base is built upon the DeepSeek-Coder-Base-v1.5 7B model \citep{deepseek-coder}, as our analysis indicates that initializing from a code-pretrained model provides superior results compared to starting with a general-purpose LLM. Moreover, we observe that mathematical pre-training not only enhances mathematical capabilities but also yields improvements on the MMLU \citep{mmlu} and BBH \citep{bbh} benchmarks, suggesting a boost in general reasoning performance beyond just mathematical tasks.

Following the pre-training stage, DeepSeekMath-Base undergoes further refinement through mathematical instruction tuning, incorporating chain-of-thought \citep{cot}, program-of-thought \citep{pot,pal}, and tool-integrated reasoning \citep{tora} datasets. The finalized model, DeepSeekMath-Instruct 7B, surpasses all other 7B models and demonstrates competitiveness with instruction-tuned models at the 70B parameter scale within the open-source community.

Additionally, we present Group Relative Policy Optimization (GRPO), a novel reinforcement learning (RL) approach derived from Proximal Policy Optimization (PPO) \citep{schulman2017proximal}. GRPO eliminates the need for a critic model by using group-based scoring to estimate the baseline, thereby offering significant reductions in computational cost. Using only a selected subset of English instruction tuning datasets, GRPO achieves marked improvements over the robust DeepSeekMath-Instruct baseline on both in-domain (GSM8K: 82.9\% $\rightarrow$ 88.2\%, MATH: 46.8\% $\rightarrow$ 51.7\%) and out-of-domain math tasks (e.g., CMATH: 84.6\% $\rightarrow$ 88.8\%) during RL training. We further outline a comprehensive framework that contextualizes diverse approaches such as Rejection Sampling Fine-Tuning (RFT) \citep{yuan2023scaling}, Direct Preference Optimization (DPO) \citep{dpo}, PPO, and GRPO, revealing that each method can be understood as either a direct or an approximated RL strategy. We perform extensive empirical studies—comparing, for example, online versus offline optimization, outcome versus process supervision, and single-step versus iterative RL—to probe critical aspects of this unified approach. Lastly, we elucidate the underlying mechanisms through which RL enhances instruction-tuned model performance, and discuss future directions to further advance RL efficacy within this integrated framework.

\subsection{Contributions}
This work presents advancements in scalable mathematical pre-training, in addition to an in-depth investigation and assessment of reinforcement learning approaches.

\textbf{Math Pre-Training at Scale}

\begin{itemize}[topsep=0pt]
    \item We demonstrate that Common Crawl, an openly available data resource, offers a rich foundation for mining mathematical content. By employing a carefully engineered data curation pipeline, we assemble the DeepSeekMath Corpus—a refined, math-specific dataset comprising 120B tokens extracted from web documents—outstripping the math web page datasets used by Minerva \citep{minerva} by a factor of nearly 7, and exceeding the size of OpenWebMath \citep{openwebmath} by approximately 9 times.

    \item Our DeepSeekMath-Base 7B model, upon pre-training, reaches performance levels comparable to those of Minerva 540B \citep{minerva}. This suggests that model scale alone does not solely determine mathematical reasoning proficiency; smaller models pre-trained with high-caliber data can likewise achieve substantial competence.

    \item The outcomes of our mathematical training experiments reveal that pre-training on code prior to math pre-training notably augments models’ problem-solving proficiency, both when external tools are available and when they are not. This provides partial resolution to the persistent query: \textit{does code training enhance reasoning capability?} Our evidence supports that, at least regarding mathematical reasoning, it does.

    \item Despite its frequent inclusion, particularly in math-focused research, the use of arXiv paper datasets in training yields no significant improvement across the suite of mathematical benchmarks utilized in our evaluations.
\end{itemize}

\textbf{Exploration and Analysis of Reinforcement Learning}

\begin{itemize}[topsep=0pt]
    \item We present Group Relative Policy Optimization (GRPO), a reinforcement learning algorithm that prioritizes efficiency and efficacy. Unlike approaches reliant on a critic model, GRPO utilizes group scores to approximate the baseline, thereby markedly diminishing the computational burden typically associated with methods such as Proximal Policy Optimization (PPO).
    \item We show that the application of GRPO leads to substantial improvements in the performance of our instruction-tuned DeepSeekMath-Instruct model, exclusively utilizing instruction-tuning datasets. Additionally, we find that this approach also improves generalization to out-of-domain tasks during reinforcement learning.
    \item We propose a cohesive framework for understanding various approaches—including RFT, DPO, PPO, and GRPO. A broad array of experiments is conducted, covering online versus offline training, outcome-based versus process-based supervision, as well as comparisons between single-turn and iterative reinforcement learning, providing in-depth insights into key aspects of this framework.
    \item Building upon our unified perspective, we delve into the factors underpinning reinforcement learning effectiveness and highlight several promising directions for enhancing the reinforcement learning of large language models (LLMs).
\end{itemize}

\subsection{Summary of Evaluations and Metrics}

\begin{itemize}[topsep=0pt]
    \item \textbf{English and Chinese Mathematical Reasoning}: We rigorously evaluate our models on both English and Chinese mathematical datasets spanning grade-school to university-level problems. The English evaluations draw on datasets such as GSM8K \citep{gsm8k}, MATH \citep{MATH}, SAT \citep{llemma}, OCW Courses \citep{minerva}, and MMLU-STEM \citep{mmlu}, while the Chinese benchmarks include MGSM-zh \citep{mgsm}, CMATH \citep{wei2023cmath}, Gaokao-MathCloze \citep{agieval}, and Gaokao-MathQA \citep{agieval}. We assess not only the capability of models to produce comprehensive text-based solutions independently of external tools but also their effectiveness in resolving mathematical problems through Python execution.

    On English datasets, DeepSeekMath-Base performs on par with the proprietary Minerva 540B model \citep{minerva}, and clearly outperforms all available open-source base models—including Mistral 7B \citep{mistral} and Llemma-34B \citep{llemma}—regardless of whether they underwent mathematical pre-training. The advantage of DeepSeekMath-Base is even more pronounced in Chinese evaluations, a result we attribute to our decision to incorporate both high-quality non-English and English mathematical corpora, diverging from previous studies \citep{minerva,llemma} that focused exclusively on English pre-training. By further applying mathematical instruction tuning and reinforcement learning, DeepSeekMath-Instruct and DeepSeekMath-RL yield robust results, notably achieving over 50\% accuracy on the competition-grade MATH dataset, marking a first among open-source efforts.
\end{itemize}

\item \textbf{Formal Mathematics}: We assess DeepSeekMath-Base through the informal-to-formal theorem proving benchmark detailed in \citep{dsp_proof} on the miniF2F dataset \citep{minif2f}, using Isabelle \citep{isabelle} as the designated proof assistant. DeepSeekMath-Base exhibits robust few-shot capabilities for automatic formalization.

\item \textbf{Natural Language Understanding, Reasoning, and Code}: To develop a thorough evaluation of the model’s abilities in understanding language, reasoning, and programming, we test DeepSeekMath-Base on the Massive Multitask Language Understanding (MMLU) benchmark \citep{mmlu}, featuring 57 multiple-choice subject areas, and on BIG-Bench Hard (BBH) \citep{bbh}, which includes 23 demanding multi-step reasoning problems. In addition, HumanEval \citep{codex} and MBPP \citep{mbpp} are utilized as standard code generation benchmarks. We observe that math-centric pre-training improves both comprehension and reasoning tasks.

\section{Math Pre-Training}

\subsection{Data Collection and Decontamination} Here, we describe our methodology for building the DeepSeekMath Corpus using Common Crawl data. Figure \ref{fig:data_collect} illustrates the iterative strategy for systematically extracting a substantial mathematical dataset, initiated from a seed corpus comprised of high-quality mathematics resources. This iterative framework can be adapted for additional domains like programming.

Our initial seed corpus is OpenWebMath \citep{openwebmath}, recognized for its curated mathematical text quality. With this foundation, we train a fastText model \citep{joulin2016fasttext} to identify further mathematical webpages similar to those in OpenWebMath. The training process involves 500,000 randomly sampled seed documents labeled as positive examples and 500,000 randomly sampled Common Crawl pages as negatives. The open-source implementation\footnote{\url{https://fasttext.cc}} is employed, using 256-dimensional vectors, a learning rate of 0.1, word n-grams of up to length 3, a minimum word frequency of 3, and 3 epochs. To manage the scale of Common Crawl, URL-based deduplication along with near-duplicate removal is performed, reducing the corpus to 40B HTML documents. The trained fastText classifier is then used to identify and retrieve math-relevant pages. To ensure content quality, retrieved pages are ranked based on fastText scores, retaining only those with the highest predicted relevance. Subsequent data retention thresholds are established through preliminary pre-training runs evaluating subsets at 40B, 80B, 120B, and 160B tokens, with the first iteration maintaining the top 40B tokens.

Following the initial round of data collection, a significant number of mathematical web pages are still missing from the dataset. This shortfall primarily arises because the fastText model was originally trained on a set of positive samples that does not exhibit enough variety. To address this, we seek out further mathematical web domains to broaden and diversify our seed corpus, thereby refining the fastText model. Our process begins by partitioning the full Common Crawl dataset into non-overlapping domains, where a domain corresponds to web pages under a common base URL. For each domain, we determine the proportion of pages that were successfully harvested during the first iteration. If a domain has more than 10\% of its web pages collected, we classify it as math-related (such as \url{mathoverflow.net}).

Next, we manually review and label URLs within these math-related domains that correspond to mathematical material (for example, \url{mathoverflow.net/questions}). Any web pages associated with these URLs that were missed in the earlier round are incorporated into our seed corpus. This strategy increases the quantity of positive training examples, enabling the subsequent training of a fastText model with enhanced capability to retrieve mathematical data in future iterations. After four cycles of collection, we acquire 35.5M mathematical web pages, which aggregate to 120B tokens. Notably, by the fourth round, nearly 98\% of all content had already been collected in the third, leading us to discontinue further data collection at this stage.

In order to prevent overlap with evaluation benchmarks, we implement a filtering process as described in \cite{deepseek-coder}. Specifically, web pages are removed from our math training corpus if they contain questions or answers drawn from English math benchmarks like GSM8K \citep{gsm8k} and MATH \citep{MATH}, as well as Chinese resources such as CMATH \citep{wei2023cmath} and AGIEval \citep{agieval}. The removal is conducted as follows: if any segment of text includes an exact match for a 10-gram sequence from the evaluation benchmarks, it is excluded. For benchmark strings shorter than 10 grams but consisting of at least 3 grams, we also apply exact matching and filter accordingly.

\subsection{Assessing the Quality of the DeepSeekMath~Corpus}

To assess the value of the DeepSeekMath~Corpus relative to other contemporary mathematical corpora, we conduct pre-training experiments comparing it to recent math-focused datasets:

\begin{itemize}[topsep=0pt]
    \item \textbf{MathPile} \citep{mathpile}: a multi-source collection (8.9B tokens) drawing from textbooks, Wikipedia, ProofWiki, CommonCrawl, StackExchange, and arXiv, with over 85\% of its contents derived from arXiv;
    \item \textbf{OpenWebMath} \citep{openwebmath}: curated from CommonCrawl data by filtering for mathematical content, comprising 13.6B tokens;
    \item \textbf{Proof-Pile-2} \citep{llemma}: a corpus that merges OpenWebMath, AlgebraicStack (10.3B tokens of mathematical code), and arXiv articles (28.0B tokens). For Proof-Pile-2 experiments, we adopt the arXiv:Web:Code ratio of 2:4:1 following the methodology in \cite{llemma}.
\end{itemize}

\subsubsection{Training Setting}
\label{sec:quality-policy}
We conduct math-specific training on a general-purpose language model with 1.3B parameters, which adheres to the same architecture as DeepSeek LLMs \citep{deepseek-llm} and is referred to here as DeepSeek-LLM 1.3B. Each mathematical corpus is used to train a distinct model for 150B tokens. All training is performed using the streamlined and efficient HAI-LLM \citep{haillm} infrastructure. Consistent with the methodology employed for DeepSeek LLMs, optimization utilizes AdamW \citep{adamW} with settings $\beta_1=0.9$, $\beta_2=0.95$, and $\mathrm{weight\_decay}=0.1$. The learning rate follows a multi-phase schedule: it peaks after 2,000 warmup steps, declines to 31.6\% of its maximum after 80\% of training steps, and further drops to 10.0\% of the maximum by the 90\% mark. The peak learning rate is set at 5.3e-4. We adopt a batch size of 4M tokens and a context window of 4K tokens.

\subsubsection{Evaluation Results}

\textbf{DeepSeekMath~Corpus stands out for its high quality, multilingual scope, and unprecedented scale.}

\begin{itemize}[topsep=0pt]
    \item \textbf{High-quality}:
    To measure downstream effectiveness, we assess few-shot chain-of-thought prompting \cite{cot} across 8 mathematical benchmark datasets. The results, displayed in Table \ref{tab:corpora_comparison}, indicate that models trained on the DeepSeekMath~Corpus consistently outperform those trained on alternatives. Furthermore, as illustrated in Figure \ref{fig:corpora_comparisons}, the DeepSeekMath-trained model surpasses the Proof-Pile-2-trained model even after processing 50B tokens (equivalent to one epoch over Proof-Pile-2), signaling superior average quality within DeepSeekMath~Corpus.
    \item \textbf{Multilingual}:
    The DeepSeekMath~Corpus contains materials spanning numerous languages, with English and Chinese being most prevalent. Table \ref{tab:corpora_comparison} demonstrates that models trained on DeepSeekMath~Corpus see improved mathematical reasoning in both English and Chinese. By comparison, the predominately English-oriented corpora used in prior work offer minimal, and sometimes detrimental, effects on mathematical reasoning in Chinese.
    \item \textbf{Large-scale}:
    DeepSeekMath~Corpus vastly exceeds previous mathematical corpora in size. Figure \ref{fig:corpora_comparisons} reveals that DeepSeek-LLM 1.3B, after training on DeepSeekMath~Corpus, not only learns more rapidly but also sustains performance improvements over time. On the other hand, baseline datasets are notably smaller, often repeated throughout training, leading to early saturation in model capabilities.
\end{itemize}

\subsection{Training and Evaluating DeepSeekMath-Base 7B}

This section details the development and evaluation of DeepSeekMath-Base 7B, a foundational model designed for advanced mathematical reasoning. The model is built upon DeepSeek-Coder-Base-v1.5 7B \citep{deepseek-coder} and is trained using 500B tokens. The dataset comprises 56\% from the DeepSeekMath Corpus, 4\% from AlgebraicStack, 10\% from arXiv, 20\% from Github code repositories, and the final 10\% consists of English and Chinese natural language samples sourced from Common Crawl. We generally adhere to the training protocol outlined in Section \ref{sec:quality-policy}, with the exceptions that the peak learning rate is increased to 4.2e-4 and the batch size is configured at 10M tokens.

A thorough evaluation of DeepSeekMath-Base 7B is conducted to measure its mathematical proficiency, particularly its capability to independently construct complete solutions, tackle mathematical tasks with external tool support, and carry out rigorous theorem proving. Additionally, we assess the broader capabilities of the base model, including its aptitude in natural language comprehension, logical reasoning, and programming.

\paragraph{Step-by-Step Reasoning in Mathematical Problem Solving}

DeepSeekMath-Base’s competence in addressing mathematical questions is tested using few-shot chain-of-thought prompting \citep{cot} across eight English and Chinese benchmarks. The selected benchmarks span quantitative reasoning tasks—such as GSM8K \citep{gsm8k}, MATH \citep{MATH}, and CMATH \citep{wei2023cmath}—and multiple-choice assessments—like MMLU-STEM \citep{mmlu} and Gaokao-MathQA \citep{agieval}—thereby covering a wide array of mathematical topics from elementary arithmetic through to college-level material.

As indicated in Table \ref{tab:base_math_cot}, DeepSeekMath-Base 7B consistently achieves superior results across all eight evaluation benchmarks when compared to other open-source base models—including the commonly adopted general-purpose Mistral 7B \citep{mistral} and the newer Llemma 34B \citep{llemma}, the latter being specifically trained on the Proof-Pile-2 math corpus \citep{llemma}. Particularly striking is its performance on the MATH dataset, representative of competition-level mathematical problems, where DeepSeekMath-Base 7B exceeds the top-performing open-source base models by a margin exceeding 10\% absolute. It also outperforms Minerva 540B \citep{minerva}, a proprietary model which, despite its substantially larger scale at 540B parameters and additional training on mathematics resources via PaLM \citep{palm}, trails behind DeepSeekMath-Base 7B.

\paragraph{Mathematical Problem Solving with Tool Use}  
To assess capabilities in program-assisted mathematical reasoning, we conduct few-shot program-of-thought prompting experiments \citep{pot,pal} on GSM8K and MATH datasets. In these tasks, models are instructed to solve presented problems by generating executable Python scripts, making use of modules such as \textit{math} and \textit{sympy} for complex numerical operations. The correctness of each solution is based on the output generated by the corresponding script. According to Table \ref{tab:base_math_pot_proof}, DeepSeekMath-Base 7B achieves performance surpassing that of Llemma 34B, which was previously the strongest open-source model.

\paragraph{Formal Mathematics}

Automating formal proof construction offers substantial advantages by improving the rigor and reliability of mathematical arguments while increasing proof generation efficiency—an area garnering increasing research focus. For this purpose, we measure the informal-to-formal translation capability of DeepSeekMath-Base 7B using the methodology of \citep{dsp_proof}, which requires transforming an informal conjecture and proof into a formal proof, guided by both the informal and formal problem statements. Evaluations on the miniF2F benchmark \citep{minif2f}—which tests Olympiad-level problems—employ few-shot prompting to generate Isabelle proofs for each instance. Adopting the process in \cite{dsp_proof}, we prompt models to produce high-level proof sketches that are then completed by Sledgehammer \citep{sledgehammer}, an automated prover tool, to address lower-level proof obligations. Table \ref{tab:base_math_pot_proof} reveals that DeepSeekMath-Base 7B attains strong results in automated formalization tasks.

\paragraph{Natural Language Understanding, Reasoning, and Code}

To gauge broader capabilities, we evaluate DeepSeekMath-Base 7B on tasks assessing natural language comprehension (MMLU \citep{mmlu}), logical reasoning (BBH \citep{bbh}), and programming proficiency (HumanEval \citep{codex} and MBPP \citep{mbpp}). Table \ref{tab:understanding_reasoning_code} illustrates that, compared to its predecessor DeepSeek-Coder-Base-v1.5 \citep{deepseek-coder}, DeepSeekMath-Base 7B achieves notable advancements on MMLU and BBH, highlighting the beneficial effect of focused mathematical training on language understanding and reasoning abilities. Moreover, by incorporating additional code tokens during continued training, DeepSeekMath-Base 7B is able to sustain the performance levels of DeepSeek-Coder-Base-v1.5 on coding tasks. Collectively, DeepSeekMath-Base 7B markedly outperforms the general-purpose Mistral 7B \citep{mistral} across all three reasoning and coding benchmarks.

\section{Supervised Fine-Tuning}

\subsection{SFT Data Curation}

We assembled a dataset tailored for instruction-tuning in mathematics, which comprises problems presented in both English and Chinese. These problems originate from a variety of mathematical domains and span a spectrum of difficulties. Each problem is matched with a detailed solution utilizing chain-of-thought (CoT) \citep{cot}, program-of-thought (PoT) \citep{pot,pal}, or tool-integrated reasoning \citep{tora} approaches. The dataset includes a total of 776,000 training pairs.

\begin{itemize}[topsep=0pt]
    \item \textbf{English mathematical datasets}: We have provided tool-integrated annotations for both GSM8K and MATH problems, and incorporated selected items from MathInstruct \citep{MathInstruct} as well as the training partition of Lila-OOD \citep{lila}, with solutions generated through CoT and PoT methodologies. The English dataset encompasses numerous mathematical subfields, such as algebra, probability, number theory, calculus, and geometry.
    \item \textbf{Chinese mathematical datasets}: Our compilation consists of K-12 mathematics problems in Chinese, covering 76 distinct topics, including but not limited to linear equations. Each item is furnished with both CoT and tool-integrated reasoning solutions.
\end{itemize}

\subsection{Training and Evaluating DeepSeekMath-Instruct 7B}

This subsection describes the development of DeepSeekMath-Instruct 7B, which is fine-tuned for mathematical instruction using DeepSeekMath-Base as its foundation. For training, examples are concatenated at random, ensuring the total context does not surpass 4K tokens per input. The model is optimized for 500 iterations using a batch size of 256 and a fixed learning rate of $5 \times 10^{-5}$.

To assess the mathematical reasoning abilities of the models—with and without external tool usage—we use four quantitative reasoning benchmarks in both English and Chinese. We compare our system with leading contemporaneous models:

\begin{itemize}[topsep=0pt]
    \item \textbf{Closed-source models}: These include (1) models from the GPT series, such as GPT-4 \citep{gpt4} and GPT-4 Code Interpreter~\footnote{\url{https://openai.com/blog/chatgpt-plugins\#\#code-interpreter}}, recognized as the strongest representatives, (2) Gemini Ultra and Pro \citep{gemini}, (3) Inflection-2 \citep{inflection-2}, (4) Grok-1~\footnote{\url{https://x.ai/model-card}}, as well as solutions recently launched by Chinese technology companies, including (5) Baichuan-3~\footnote{\url{https://www.baichuan-ai.com}} and (6) the latest model from the GLM family, GLM-4~\footnote{\url{https://open.bigmodel.cn/dev/api\#glm-4}} \citep{glm}.
\end{itemize}

These models serve broad objectives and have generally experienced extensive alignment steps.

\item \textbf{Open-source models} encompass: foundational large language models such as (1) DeepSeek-LLM-Chat 67B \citep{deepseek-llm}, (2) Qwen 72B \citep{qwen}, (3) SeaLLM-v2 7B \citep{seallm}, and (4) ChatGLM3 6B \citep{chatglm3}. Additionally, several models are optimized for mathematical competence, including (5) InternLM2-Math 20B~\footnote{\url{https://github.com/InternLM/InternLM-Math}}, which extends InternLM2 and incorporates dedicated mathematical pretraining as well as instruction tuning, (6) Math-Shepherd-Mistral 7B, created by applying PPO training \citep{schulman2017proximal} with a process-supervised reward model to Mistral 7B \citep{mistral}, (7) the WizardMath collection \citep{wizardmath}, which advances the mathematical reasoning ability of Mistral 7B and Llama-2 70B \citep{llama2} via evolve-instruct (a version of instruction tuning using automatically generated instructions) and PPO on primarily GSM8K and MATH tasks, (8) MetaMath 70B \citep{metamath}, consisting of Llama-2 70B further trained with a supplemented GSM8K and MATH dataset, (9) ToRA 34B \cite{tora}, a variant of CodeLlama 34B that is tuned for mathematical reasoning with tool integration, and (10) MAmmoTH 70B \citep{MathInstruct}, which is Llama-2 70B tuned on the MathInstruct corpus.

As detailed in Table \ref{tab:sft_rl_math}, when restricting evaluations to settings where external tools are not permitted, DeepSeekMath-Instruct 7B achieves superior step-wise reasoning performance. In particular, for the competition-level MATH benchmark, our model exceeds all other open-source systems and surpasses most proprietary alternatives (such as Inflection-2 and Gemini Pro) by a margin of at least 9 percentage points. This result persists even when comparing with considerably larger models (e.g., Qwen 72B) and with models benefiting from specialized math-oriented reinforcement learning strategies (for instance, WizardMath-v1.1 7B). While DeepSeekMath-Instruct's accuracy on MATH matches leading Chinese proprietary systems like GLM-4 and Baichuan-3, it still does not attain the performance of GPT-4 and Gemini Ultra.

Allowing for the use of both natural language reasoning and programmatic tools in problem-solving, DeepSeekMath-Instruct 7B attains nearly 60\% accuracy on MATH, surpassing all publicly available models. Across other datasets, it performs competitively relative to DeepSeek-LLM-Chat 67B, the previous top open-source performer, despite being an order of magnitude smaller.

\section{Reinforcement Learning}

\subsection{Group Relative Policy Optimization} Reinforcement learning (RL) has consistently been shown to bolster the mathematical reasoning proficiency of LLMs following Supervised Fine-Tuning (SFT) \citep{wang2023math,wizardmath}. Here, we propose Group Relative Policy Optimization (GRPO), an RL algorithm designed to be both resource-efficient and highly effective.

\subsubsection{From PPO to GRPO} Proximal Policy Optimization (PPO) \citep{schulman2017proximal}, a standard actor-critic RL method, is widely employed for RL-based fine-tuning of LLMs \citep{ouyang2022training}. Specifically, PPO enhances LLMs by maximizing this surrogate objective function:

\begin{equation}
\footnotesize
    \mathcal{J}_{PPO}(\theta) = \mathbb{E}{[q \sim P(Q), o \sim \pi_{\theta_{old}}(O|q)]} \frac{1}{|o|} \sum_{t=1}^{|o|} \min \left[ \frac{\pi_\theta(o_{t} | q, o_{<t})}{\pi_{\theta_{old}}(o_{t} | q, o_{<t})} A_{t}, \text{clip} \left( \frac{\pi_\theta(o_{t} | q, o_{<t})}{\pi_{\theta_{old}}(o_{t} | q, o_{<t})}, 1 - \varepsilon, 1 + \varepsilon \right)

In this context, $\pi_{\theta}$ denotes the current policy, while $\pi_{\theta_{old}}$ refers to the previous version. The variables $q$ and $o$ represent questions and outputs drawn from the question set and from the output distribution generated by $\pi_{\theta_{old}}$, respectively. The scalar $\varepsilon$ acts as a hyper-parameter associated with the clipping mechanism introduced in PPO to stabilize optimization. $A_t$ signifies the advantage, which is derived using the Generalized Advantage Estimation (GAE) method \citep{gae}, relying on reward sequences $\{r_{\ge t}\}$ and an associated value function $V_{\psi}$ learned concurrently. Consequently, PPO entails training both a value function and a policy; to counteract potential overfitting to the reward model, a KL divergence penalty term between the policy under training and a reference model is typically incorporated into the per-token reward, following the standard practice from \citep{ouyang2022training}:

\begin{equation}
    r_{t} = r_\varphi(q, o_{\le t}) - \beta \log\frac{\pi_{\theta}(o_{t}|q, o_{<t})}{\pi_{ref}(o_{t}|q, o_{<t})},
\label{eq:PPO-reward}
\end{equation}

Here, $r_\varphi$ corresponds to the reward model, $\pi_{ref}$ is the reference model—generally initialized from the SFT model—and $\beta$ is the scaling coefficient for the KL regularization.

Given that the value function in PPO is frequently modeled using an architecture of similar scale as the policy network, this results in significant memory usage and computational load. Furthermore, in RL applications, the value function provides a baseline for computing the advantage and reducing variance. In the case of LLMs, it is typical for only the reward associated with the final token to be available from the reward model, which complicates the accurate training of a value function for every token. To overcome these obstacles, as depicted in Figure \ref{fig:grpo}, we introduce the Group Relative Policy Optimization (GRPO) technique. GRPO dispenses with the explicit value function approximation found in PPO, instead substituting as a baseline the mean reward from a batch of sampled outputs generated for a single question. More concretely, for each question $q$, GRPO draws a set of outputs $\{o_1, o_2, \cdots, o_G\}$ via the old policy $\pi_{\theta_{old}}$ and then updates the policy by maximizing the following objective function:

\begin{equation}
\footnotesize
\begin{split}
    \mathcal{J}_{GRPO}(\theta) &= \mathbb{E}{[q \sim P(Q), \{o_i\}_{i=1}^G \sim \pi_{\theta_{old}}(O|q)]}  \\
    & \frac{1}{G}\sum_{i=1}^G\frac{1}{|o_i|} \sum_{t=1}^{|o_i|} \left\{ \min \left[ \frac{\pi_\theta(o_{i,t} | q, o_{i,<t})}{\pi_{\theta_{old}}(o_{i,t} | q, o_{i,<t})} \hat{A}_{i,t}, \text{clip} \left( \frac{\pi_\theta(o_{i,t} | q, o_{i,<t})}{\pi_{\theta_{old}}(o_{i,t} | q, o_{i,<t})}, 1 - \varepsilon, 1 + \varepsilon \right)  \hat{A}_{i,t} \right] - \beta \mathbb{D}_{KL}\left[\pi_{\theta} || \pi_{ref}\right]\right\} ,
\end{split}
\label{eq:GRPO-obj}
\end{equation}

In this formulation, $\varepsilon$ and $\beta$ remain as hyper-parameters, and $\hat{A}_{i,t}$ stands for the advantage computed by contrasting the rewards of outputs sampled in the same group—details of which are elaborated upon in subsequent sections. By adopting a group-relative approach to advantage computation, GRPO is naturally suited to exploit the pairwise structure inherent in many reward modeling datasets, which are based on comparing responses to identical prompts. Importantly, unlike methods that inject the KL penalty directly into the reward, GRPO incorporates KL divergence regularization between the policy being trained and the reference policy straight into the objective, thus streamlining the calculation of $\hat{A}_{i,t}$. Notably, whereas the

\begin{algorithm}[t]
  \small
  \caption{Iterative Group Relative Policy Optimization}
  \textbf{Input} Initial policy model $\pi_{\theta_{\text{init}}}$; reward models $r_\varphi$; task prompt set $\mathcal{D}$;
  hyperparameters $\varepsilon$, $\beta$, $\mu$
  \begin{algorithmic}[1]
    \State Assign policy model: $\pi_\theta \leftarrow \pi_{\theta_{\text{init}}}$
    \For{iteration = 1, \dots, I}
       \State Set reference model: $\pi_{ref} \leftarrow \pi_{\theta}$
      \For{step = 1, \dots, M}
      \State Randomly select batch $\mathcal{D}_b$ from $\mathcal{D}$
      \State Update previous policy: $\pi_{\theta_{old}} \leftarrow \pi_{\theta}$
      \State For every question $q$ in $\mathcal{D}_b$, sample $G$ outputs $\{o_i\}_{i=1}^G \sim \pi_{\theta_{old}} (\cdot \mid q) $
      \State Evaluate each output $o_i$ to get rewards $\{r_i\}_{i=1}^{G}$ using $r_{\varphi}$
      \State Employ group relative advantage estimation to compute $\hat{A}_{i,t}$ for the $t$-th token in each $o_i$
      \For{GRPO iteration = 1, \dots, $\mu$}
        \State Adjust policy model $\pi_{\theta}$ by maximizing the GRPO objective from Equation \ref{eq:GC-GRPO}
      \EndFor
    \EndFor \State Continuously improve $r_\varphi$ via replay-based training.
    \EndFor
  \end{algorithmic}
  \textbf{Output} $\pi_\theta$
  \label{alg:iter-grpo}
\end{algorithm}

\subsubsection{Outcome Supervision RL with GRPO} To define the method precisely, consider a set of responses $\{o_1, o_2, \cdots, o_G\}$ sampled for each question $q$ from the former policy $\pi_{\theta_{old}}$. Each output is then evaluated with a reward model, producing a set of $G$ reward values $\mathbf{r}=\{r_1, r_2, \cdots, r_G\}$. The reward values are standardized by subtracting the group mean and dividing by the group standard deviation. Under outcome supervision, each output $o_i$ is assigned a normalized reward at its completion, and this value is uniformly attributed as the advantage $\hat{A}_{i, t}$ for every token within $o_i$: that is, $\hat{A}_{i, t} = \widetilde{r}_i = \frac{r_i- {\rm mean}(\mathbf{r})}{{\rm std}(\mathbf{r})}$. The policy is then updated by maximizing the objective provided in equation (\ref{eq:GRPO-obj}).

\subsubsection{Process Supervision RL with GRPO} Relying solely on outcome rewards may be limiting for sophisticated mathematical problems, as outcome supervision restricts feedback to only the terminal point of each output. Inspired by \cite{wang2023math}, we extend our investigation to process supervision, which furnishes reward feedback after every reasoning stage. Specifically, for question $q$ and $G$ sampled outputs $\{o_1, o_2, \cdots, o_G\}$, a process reward model evaluates each intermediate step, producing structured rewards: $\mathbf{R} = \{ \{r_1^{index(1)},\cdots,r_1^{index(K_1)}\}, \cdots,  \{r_G^{index(1)},\cdots,r_G^{index(K_G)}\} \}$, where $index(j)$ identifies the position of the last token in the $j$-th step and $K_i$ denotes the number of steps in output $i$. These step-level rewards are also normalized, calculated as $\widetilde{r}_i^{index(j)} = \frac{r_i^{index(j)} - {\rm mean(\mathbf{R})}}{{\rm std(\mathbf{R})}}$. The token-level advantage $\hat{A}_{i, t}$ is then assigned as the sum of the normalized rewards for all subsequent (and current) steps, that is, $\hat{A}_{i, t} = \sum_{index(j) \ge t} \widetilde{r}_i

\subsubsection{Iterative RL with GRPO} During the reinforcement learning procedure, reliance on an earlier reward model may become inadequate for guiding the present policy model. Consequently, we also examine iterative RL leveraging GRPO. As depicted in Algorithm \ref{alg:iter-grpo}, iterative GRPO entails periodically constructing new reward model training datasets from fresh samples drawn by the policy model, while continuing to update the previous reward model using a replay buffer that includes 10\% of previously seen data. Subsequently, we designate the policy model as the reference model, and further refine the policy model utilizing the updated reward model.

\subsection{Training and Evaluating DeepSeekMath-RL}

Our reinforcement learning experiments are carried out using DeepSeekMath-Instruct 7B as the policy model. The RL training data comprises chain-of-thought-format questions associated with GSM8K and MATH, sourced from the SFT dataset, amounting to approximately 144,000 questions. We omit other SFT questions to precisely assess the influence of RL on evaluation sets that do not have coverage during the RL phase. Construction of reward model training data follows the protocol from \citep{wang2023math}. The reward model is initially trained starting from DeepSeekMath-Base 7B, employing a learning rate of 2e-5. In GRPO, we set the policy model learning rate to 1e-6 and adopt a KL penalty coefficient of 0.04. For each question, 64 candidate responses are sampled. The maximum sequence length is set to 1024, and training batches also consist of 1024 samples. The policy model undergoes a single update after each round of exploration. DeepSeekMath-RL 7B is assessed on evaluation benchmarks in the same manner as DeepSeekMath-Instruct 7B. For DeepSeekMath-RL 7B, GSM8K and MATH with chain-of-thought reasoning formats are treated as in-domain benchmarks, while the rest are categorized as out-of-domain.

Table \ref{tab:sft_rl_math} presents the results achieved by both open- and closed-source models, using chain-of-thought and tool-augmented reasoning methods on English and Chinese evaluation sets. Our key observations are: 1) DeepSeekMath-RL 7B achieves accuracy rates of 88.2\% on GSM8K and 51.7\% on MATH by leveraging chain-of-thought reasoning. This positions it ahead of every open-source model within the 7B to 70B parameter range and places it above most closed-source models as well. 2) Notably, DeepSeekMath-RL 7B's training relies solely on chain-of-thought-structured instruction data from GSM8K and MATH, building upon DeepSeekMath-Instruct 7B. Even with this limited training distribution, it consistently outperforms DeepSeekMath-Instruct 7B on all tested metrics, underscoring the impact of reinforcement learning.

\section{Discussion}
This section summarizes our insights gleaned from both pre-training and RL-based experiments.

\subsection{Lessons Learnt in Pre-Training}

We begin by outlining the main takeaways from our pre-training procedures. Unless otherwise noted, we maintain the experimental configuration described in Section \ref{sec:quality-policy}. It should be mentioned that, throughout this section, "DeepSeekMath Corpus" specifically refers to the 89B-token dataset produced in the data collection process's second cycle.

\subsubsection{Code Training Benefits Mathematical Reasoning}

There is a widely held, though often untested, assumption that training on code enhances reasoning abilities. Here, we provide partial empirical support for this hypothesis, particularly as it relates to mathematical reasoning: incorporating code into training demonstrably enhances a model’s capacity for mathematical reasoning, both with and without supplementary tools.

To probe the effect of code training on mathematical reasoning abilities, we compared two different pre-training paradigms: a sequential two-stage training regimen and a unified one-stage training approach:

\textbf{Two-Stage Training}

\begin{itemize}[topsep=0pt]
    \item \textbf{Code Training for 400B Tokens $\rightarrow$ Math Training for 150B Tokens}: DeepSeek-LLM 1.3B undergoes initial pretraining on 400B tokens of code, which is subsequently followed by an additional 150B tokens focused on mathematical content.
    \item \textbf{General Training for 400B Tokens $\rightarrow$ Math Training for 150B Tokens}: For a comparative baseline, the same procedure is carried out using 400B general domain tokens—sourced from DeepSeek-AI’s large general-purpose dataset—rather than code tokens in the first phase, enabling an analysis of the potential benefits of code-centric pretraining for downstream mathematical reasoning abilities.
\end{itemize}

\textbf{One-Stage Training}

\begin{itemize}[topsep=0pt]
    \item \textbf{Math Training for 150B Tokens}: The DeepSeek-LLM 1.3B model is directly trained using 150B mathematical tokens only.
    \item \textbf{Training on a mixture of 400B Code Tokens and 150B Math Tokens}: Observing that sequential math training following code pretraining negatively affects coding competency, we assess whether a combined (simultaneous) training on 400B code and 150B math tokens sustains gains in mathematical reasoning while reducing the detrimental effects of catastrophic forgetting.
\end{itemize}

\paragraph{Results}
The empirical results from Table \ref{tab:code-math} and Table \ref{tab:code-math-general-eval} illustrate how downstream tasks are impacted by the varying training methodologies.

Incorporating code tokens in pretraining shows marked improvement in program-assisted mathematical reasoning under both sequential and joint training strategies. Specifically, Table \ref{tab:code-math} demonstrates that even standalone code pretraining greatly boosts the model’s proficiency at tasks like GSM8K and MATH (in Python-based settings). Further mathematical finetuning after code training delivers additional gains. Notably, simultaneous training with a mixture of code and math tokens not only lessens catastrophic forgetting—which is observed during the two-step approach—but also offers compounded improvements in both coding capabilities (Table \ref{tab:code-math-general-eval}) and program-based mathematical reasoning (Table \ref{tab:code-math}).

Code-focused pretraining also has beneficial effects on mathematical reasoning without relying on external tools. In the two-step training protocol, the initial code-focused stage provides a meaningful increase in baseline performance and augments the effectiveness of the subsequent math-specific training, culminating in optimal outcomes. In contrast, training on a mixture of code and math tokens within a single stage tends to diminish math reasoning proficiency without external tool usage. A possible explanation is that the parameter capacity of DeepSeek-LLM 1.3B is insufficient to concurrently master both code and mathematical tasks when they are presented together.

\subsubsection{ArXiv Papers Show Limited Effectiveness for Mathematical Reasoning Enhancement}

The inclusion of ArXiv papers as a major source in mathematical pre-training datasets is prevalent \citep{minerva,gpt-f,llemma,mathpile}. Nevertheless, a thorough investigation of their actual influence on mathematical reasoning abilities remains largely unexplored. Contrary to common assumptions, our empirical findings suggest that arXiv papers offer little to no benefit in fostering mathematical reasoning. We evaluated models at varying scales, notably DeepSeek-LLM 1.3B and DeepSeek-Coder-Base-v1.5 7B \citep{deepseek-coder}, each trained on differently processed arXiv datasets:

\begin{itemize}[topsep=0pt]
    \item \textbf{MathPile} \citep{mathpile}: a corpus containing 8.9B tokens, curated with substantial cleaning and heuristic filtering, in which over 85\% of the tokens are sourced from scientific arXiv publications;
    \item \textbf{ArXiv-RedPajama} \citep{redpajama}: a dataset comprising the complete set of arXiv LaTeX files, post-processing to remove preambles, comments, macros, and bibliographic material, totaling 28.0B tokens.
\end{itemize}

For our analysis, we independently trained DeepSeek-LLM 1.3B on 150B tokens and DeepSeek-Coder-Base-v1.5 7B on 40B tokens from each arXiv-derived dataset. Across multiple mathematical benchmarks spanning varying difficulty levels, models exposed solely to arXiv material showed no significant progress, and in some cases, even regressions in performance. Evaluated benchmarks encompassed quantitative reasoning suites such as GSM8K and MATH (see Table \ref{tab:arxiv-cot}), multiple-choice assessments like MMLU-STEM (Table \ref{tab:arxiv-cot}), and domains of formal mathematics, e.g., miniF2F (refer to Table \ref{tab:arxiv_atp}).

It is important to note, however, that this observation is bounded by several caveats. Our current investigation does not address:

\begin{itemize}[topsep=0pt]
    \item The role of arXiv data in other mathematical tasks not included in our experiments, for example, the informalization process—transforming formalized statements or proofs into informal expressions;
    \item How arXiv-derived tokens might interact synergistically with alternative training corpora;
    \item The possibility that larger-scale models might unlock latent benefits from arXiv content.
\end{itemize}

Therefore, further detailed studies are warranted, which we identify as directions for subsequent research.

\subsection{Reinforcement Learning: Emerging Insights}
\subsubsection{Toward a Unified Framework} This section introduces a generalized analytical framework that captures various training strategies, including SFT, RFT, DPO, PPO, and GRPO, and further investigates the contributing factors within this unified formulation. The parameter gradient for such methods generally takes the following form:

\begin{equation}
    \nabla_{\theta}\mathcal{J}_{\textcolor{red}{\mathcal{A}}}(\theta) = \mathbb{E}[\underbrace{(q,o) \sim \textcolor{red}{\mathcal{D}}}_{Data \ Source}]\left( \frac{1}{|o|} \sum_{t=1}^{|o|}  \underbrace{GC_{{\mathcal{A}}}(q, o, t, \textcolor{red}{\pi_{{rf}}})}_{Gradient \ Coefficient}  \nabla_{\theta}\log \pi_{\theta}(o_t | q, o_{<t})\right).
\label{eq:objective}
\end{equation}

This objective incorporates three principal elements: 1) \textit{Data Source $\mathcal{D}$}, specifying the data utilized during training; 2) \textit{Reward Function $\pi_{{rf}}$}, responsible for providing the feedback or reward used in training; and 3) \textit{Algorithm $\mathcal{A}$}, which processes the input data alongside the reward to yield the gradient coefficient $GC$, modulating the strength of reward or penalty for each data point. Using this generalized formulation, several prominent approaches can be compared:

\begin{itemize}
    \item  \textbf{Supervised Fine-tuning (SFT)}: This method involves adjusting a pretrained model by training it further on curated human-annotated datasets.
    \item \textbf{Rejection Sampling Fine-tuning (RFT)}: RFT entails additional training of the SFT-tuned model, utilizing only the model-generated responses (in answer to SFT questions) that pass correctness-based filtering.
    \item \textbf{Direct Preference Optimization (DPO)}: DPO further improves the SFT model by fine-tuning it on a set of augmented responses, leveraging a pairwise loss to encode preference.
    \item \textbf{Online Rejection Sampling Fine-tuning (Online RFT)}: In contrast to RFT, this technique starts from the SFT model but iteratively refines the policy with augmented outputs dynamically sampled from the latest version of the policy itself.
    \item \textbf{PPO/GRPO}: These strategies initialize with the SFT model, then enhance the policy through reinforcement, training with outputs continuously drawn from the current policy.
\end{itemize}

A summary of the elements comprising these approaches can be found in Table \ref{tab:data_policy}. For an in-depth derivation, see Appendix \ref{app:analysis-rl}.

\paragraph{Insight on Data Source} We classify data sources into two main types: online sampling and offline sampling. Online sampling refers to training data obtained from the ongoing exploration behavior of the policy during training, whereas offline sampling involves data generated from the initial SFT model prior to policy updates. Both RFT and DPO operate using offline data, while Online RFT and GRPO utilize online sampling.

Figure \ref{fig:rl-analysis} illustrates that Online RFT achieves notably superior results compared to RFT across both benchmarks. During the early training phase, Online RFT performs similarly to RFT; however, as training progresses, Online RFT achieves a marked improvement, which underlines the benefits of online sampling. This can be attributed to the high similarity between the actor and SFT model initially, where differences in sampled data are negligible. In later stages, as the actor diverges from the SFT model, the actor’s sampled data contains more substantive differences, making real-time sampling increasingly advantageous.

\paragraph{Insight on Gradient Coefficient} The transformation of reward signals into gradient coefficients is fundamental for parameter updates. In our setup, the reward mechanism is categorized as either `Rule' or `Model'. The `Rule' category evaluates responses by answer accuracy, while the `Model' type involves training a reward model (based on rule-based annotations) to assign a score to each response. The critical distinction, as delineated by Equations \ref{eq:GC-RFT} and \ref{eq:GC-GRPO}, is that GRPO modifies its gradient coefficient according to the reward model’s output. This approach permits nuanced encouragement or discouragement of responses depending on the reward's value. Conversely, Online RFT treats all correct answers equally and does not penalize incorrect responses, lacking this granularity in reinforcement.

As illustrated in Figure \ref{fig:rl-analysis}, GRPO demonstrates greater effectiveness than online RFT, underscoring the advantages of adjusting both positive and negative gradient coefficients. Moreover, GRPO+PS yields better results than GRPO+OS, underscoring the merit of implementing fine-grained, step-specific gradient coefficients. We further investigate iterative RL, conducting two iterations in our experiments. Figure \ref{fig:iter-rl} indicates that iterative RL markedly enhances performance, with the most substantial improvement occurring during the initial iteration.

\subsubsection{Why RL Works?} This study employs reinforcement learning on a selected portion of instruction tuning data, achieving notable gains in model performance over the base instruction tuning model. To clarify the underlying reasons for RL's effectiveness, we assess the Pass@K and Maj@K accuracy metrics for both Instruct and RL-enhanced models across two benchmark datasets. According to Figure \ref{fig:maj-pass}, RL improves Maj@K scores, but not Pass@K. This suggests that RL primarily strengthens the robustness of the output distribution—\textbf{that is, the observed improvement stems from a higher likelihood of correct answers being present in the TopK responses, rather than an intrinsic boost to core model abilities.} Likewise, \citep{wang2023making} have identified a \textbf{misalignment problem} within SFT models in reasoning tasks, and found that a range of preference alignment techniques \citep{yuan2023rrhf,song2023preference,wang2023making} can bolster SFT models' reasoning abilities.

\subsubsection{How to Achieve More Effective RL?}
\label{sec:effec_rl}
Our results demonstrate that RL achieves strong performance in mathematical reasoning benchmarks. We also introduce a unified framework for conceptualizing various prominent training methodologies, wherein each approach is interpreted as a variant or simplification of RL. As depicted in Equation \ref{eq:objective}, this framework consists of three central elements: Data Source, Algorithm, and Reward Function. We outline prospective research opportunities pertaining to each component.

\paragraph{Data Source} The data source constitutes the foundational element for all training approaches. Specifically for RL, we define the data source as the pool of unlabeled questions and their respective responses sampled from the policy model. In this work, we restrict ourselves to questions arising from the instruction tuning process and employ straightforward nucleus sampling for response generation. This limitation could partially explain why our RL implementation yields improvements primarily for Maj@K. Looking ahead, we plan to apply our RL workflow to out-of-distribution question sets, and integrate \textbf{more sophisticated sampling (decoding) methods}—such as tree-search approaches \citep{yao2023tree}. In addition, recent advances in \textbf{efficient inference algorithms} \citep{xia-etal-2023-speculative,leviathan2023fast,kwon2023efficient,xia2024unlocking}, which govern how efficiently policy models can be explored, are likely to be of considerable significance.

\paragraph{Algorithms} Algorithms use both the data and the reward signals to determine the gradient coefficients required for updating model parameters. According to Equation \ref{eq:objective}, present approaches largely \textbf{TRUST} the reward signal, modulating the conditional probability of specific tokens up or down based solely on its direction. Nonetheless, such reliance is problematic since the reward signal is not universally dependable, particularly in highly intricate tasks. For instance, even in datasets such as PRM800K \citep{lightman2023let}, which have benefited from meticulous annotation by expert annotators, around 20\% of the annotations are erroneous\footnote{\url{https://github.com/openai/prm800k/issues/12\#issuecomment-1728491852}}. Therefore, this motivates the investigation of reinforcement learning algorithms that can withstand the effects of noisy reward signals. We anticipate that alignment techniques like \textbf{WEAK-TO-STRONG} \citep{burns2023weak} have the potential to reshape the underlying principles of learning algorithms.

\paragraph{Reward Function} The reward function serves as the primary training signal in reinforcement learning settings, typically realized as a neural reward model. There are three key research avenues concerning reward models: 1) \textbf{Enhancing generalization capabilities.} Reward models must generalize beyond their training data to address out-of-distribution scenarios and sophisticated model outputs; if not, RL may end up merely maintaining the status quo of LLM distributions instead of advancing their competence; 2) \textbf{Quantifying model uncertainty.} Incorporating measures of uncertainty could act as an interface between imprecise reward models and weak-to-strong learning frameworks; 3) \textbf{Efficient construction of high-quality process reward models} capable of delivering nuanced, step-level signals to guide the reasoning process \citep{lightman2023let,wang2023math}.

\section{Conclusion, Limitation, and Future Work} In this work, we introduce DeepSeekMath, a model that sets a new performance benchmark among open-source models on the competition-level MATH dataset, drawing near to the results achieved by proprietary systems. DeepSeekMath~is built upon DeepSeek-Coder-v1.5 7B, and has been continually trained with a corpus of 500B tokens, prominently including 120B math-related tokens curated from Common Crawl. Through comprehensive ablation analyses, we demonstrate that web pages provide substantial high-quality mathematical content, whereas arXiv contributes less than anticipated. We further propose Group Relative Policy Optimization (GRPO), an adaptation of Proximal Policy Optimization (PPO), which markedly enhances mathematical reasoning while being more memory efficient. Our empirical findings indicate that GRPO remains effective even at advanced stages of training, as evidenced by the high benchmark performance of DeepSeekMath-Instruct 7B. We also propose a unified perspective on various reinforcement learning methods, and outline prospective research paths for optimizing reinforcement learning strategies.

Despite DeepSeekMath’s notable results in quantitative reasoning, its performance in geometric reasoning and formal theorem-proving lags behind that of closed-source models. As observed in preliminary evaluations, the model struggles with geometric problems involving figures such as triangles and ellipses, which suggests possible biases in pre-training and fine-tuning data selection. Additionally, due to the model’s parameter size, DeepSeekMath~falls short of GPT-4 in few-shot learning scenarios. While GPT-4's performance significantly benefits from few-shot prompts, DeepSeekMath~displays comparable outcomes in both zero-shot and few-shot conditions. Looking forward, we plan to refine our engineered data selection processes to develop a superior pre-training corpus. Moreover, we aim to pursue further research, as discussed in Section \ref{sec:effec_rl}, to discover more efficient reinforcement learning approaches for large language models.

\end{document}